% ===================================================================
%  TABEL Nr. 7 — Het dorp in de winter. — Het landhuis in de lente.
%  Traduction 2026 d'apres texto/io, controlee sur texto/fr et
%  texto/es. Consigne : voir texto/nl/10-tabel-01.tex.
% ===================================================================

%%K t07-noto-1 noto
\VUnotes{143.2mm}{%
(*) Zie voor de bijzonderheden van de maaltijden, de tabellen 6 en 13.%
}

%%K t07-apar-1 apar
\VUsaut{4.0mm}
\VUcentre{13.2pt}{90}{TWEEDE REEKS}
\VUsaut{3.0mm}
\VUfilet{16mm}
\VUsaut{5.6mm}
\VUtitre{15.0pt}{60}{TABEL Nr. 7}
\VUsaut{3.0mm}

%%K t07-tit-1 sub
\VUcentre{12.0pt}{0}{\textit{Eerste tafereel.}}
\VUsaut{3.0mm}

%%K t07-tit-2 sub
\VUpk{10.2pt}{40}{Het dorp in de winter.}
\VUsaut{4.0mm}

%%K t07-c1-01-1 p
1. --- Tijdens de nieuwjaarsvakantie vergezelde ik mijn vader naar Nieuwstad,
een dorpje waar hij enkele handelszaken moest afhandelen.

%%K t07-c1-02-1 p
2. --- Wij namen de \VUgras{postkoets} \textsuperscript{(11)} die tussen de stad
en dat vlek rijdt; maar omdat het zeer koud was, was die reis in een weinig
comfortabel voertuig niet erg aangenaam.

%%K t07-c1-03-1 p
3. --- Daarom voelden wij een waar genoegen toen wij de \VUgras{koetsier} zijn
rijtuig op het plein zagen stilhouden, voor de \VUgras{herberg}
\textsuperscript{(2)} van de \textit{Witte Beer}. \VUgras{De herbergier}
\textsuperscript{(4)} stond ons op de drempel van zijn huis op te wachten,
rustig zijn \VUgras{pijp} rokend.

%%K t07-c1-03-2 p
Hij nodigde ons uit binnen te komen en ik liet mij niet bidden om mij voor het
takkenvuur te installeren dat in de haard knetterde. Toen lachte ik hartelijk om
de scherpe \VUgras{noordenwind} die het oude \VUgras{uithangbord}
\textsuperscript{(3)} doet knarsen en die \VUgras{de rook}
\textsuperscript{(1)} in zwarte wervelingen wegvoerde die uit de schoorsteen
kwam.

%%K t07-c1-04-1 p
4. --- Het was het uur van het middagmaal. Zonder langer te wachten aten wij met
smaak de eenvoudige maar smakelijke maaltijd die men ons opdiende (*).

%%K t07-tit-3 sub
\VUpk{10.2pt}{40}{Enkele bezoeken.}
\VUsaut{3.0mm}

%%K t07-c1-05-1 p
5. --- Na het middagmaal vergezelde ik mijn vader, die verzekeringsagent is, bij
zijn klanten. Eerst bezochten wij de \VUgras{smid} \textsuperscript{(5)}, die
naast de herberg woont. Hij was bezig een paard te beslaan. Ik bewonderde de
stevigheid van zijn maat die op zijn leren schort de \VUgras{hoef} van het paard
stevig vasthield, terwijl de smid het \VUgras{hoefijzer}
\textsuperscript{(6)} met krachtige hamerslagen bevestigde.

%%K t07-c1-06-1 p
6. --- Daarna gingen wij naar de \VUgras{schoenmaker} \textsuperscript{(7)} van
het dorp, de oude Jacob. Hoewel hij als uithangbord een prachtige \VUgras{laars}
heeft, stelde hij zich tevreden met het verzolen van een oud paar doorgesleten
schoenen.

%%K t07-c1-06-2 p
De oude man zei lachend tegen ons : « In de stad was ik “schoenmaker” en ik had
geen klanten. Hier ben ik “schoenlapper” en het werk is veel. »

%%K t07-c1-07-1 p
7. --- Hij sprak ons over zijn buurman de \VUgras{barbier}
\textsuperscript{(8)}, wiens klanten niet tevreden zijn. Zijn
\VUgras{scheermessen} zijn altijd geschaard en zijn \VUgras{scharen} knippen
nooit goed. Maar omdat hij de enige barbier van het dorp is, is men natuurlijk
gedwongen zich door hem te laten scheren en zijn haar te laten knippen. Overigens
is hij een gierig en onvriendelijk man.

%%K t07-c1-08-1 p
8. --- Gelukkig lijken de overige \VUgras{buren} niet op hem. Bijvoorbeeld de
\VUgras{wagenmaker} \textsuperscript{(9)} is een goedhartig man, vlijtig en
behulpzaam. Het is een genoegen door hem een rijtuig te laten herstellen. Zijn
werk is altijd afgewerkt.

%%K t07-c1-09-1 p
9. --- De oude Jacob klaagde ook niet over de \VUgras{zadelmaker}
\textsuperscript{(10)}, die hij soms helpt om het \VUgras{tuig} te naaien
wanneer het werk dringt.

%%K t07-c1-09-2 p
Mijn vader vroeg hem naar zijn familie : « Iedereen maakt het goed. Mijn
kleindochter zit op \VUgras{school} \textsuperscript{(12)}. Haar
\VUgras{onderwijzeres} \textsuperscript{(13)} is zeer tevreden over haar, zij is
een goede \VUgras{leerling} \textsuperscript{(14)}. Zij lijkt niet op mijn
deugniet van een zoon; hij denkt alleen aan spelen. De \VUgras{onderwijzer}
\textsuperscript{(15)} slaagt er niet in hem wat dan ook te leren. »

%%K t07-tit-4 sub
\VUsaut{3.0mm}
\VUpk{10.2pt}{40}{Dorpspraatjes.}
\VUsaut{3.0mm}

%%K t07-c1-10-1 p
10. --- De oude man vertelde ons daarna het nieuws van het dorp. Men gaat een
nieuwe school bouwen, want die welke in het \VUgras{gemeentehuis} gevestigd is
werd te klein. Overigens is dat gemakkelijk, want het zal volstaan een deel van
de tuin van de \VUgras{molenaar} te kopen. Dat zal zeer voordelig zijn voor de
arme man. Door de koude kunnen de \VUgras{molenstenen} van de \VUgras{molen}
\textsuperscript{(16)} het graan niet meer malen, want het \VUgras{rad}
\textsuperscript{(17)} werd door het \VUgras{ijs} \textsuperscript{(25)} van het
\VUgras{beekje} \textsuperscript{(23)} vastgevroren.

%%K t07-c1-11-1 p
11. --- « Wat gebouwen betreft, hebt u onze nieuwe \VUgras{kerk}
\textsuperscript{(18)} gezien? Zij is zeer schilderachtig onder haar
sneeuwmantel. De \VUgras{kraaien} \textsuperscript{(19)}, die hun oude
klokkentoren niet meer herkennen, zitten treurig op de grote boom van het
\VUgras{kerkhof} \textsuperscript{(20)} ».

%%K t07-c1-12-1 p
12. --- Die gedachte aan het kerkhof herinnerde de schoenmaker aan de mensen die
mijn vader gekend had en die vorig jaar gestorven zijn. « Hoeveel \VUgras{graven}
\textsuperscript{(21)}, mijn beste heer, zijn er bij de andere gekomen, onder de
\VUgras{cipres} \textsuperscript{(22)}! De dood heeft onze oude notaris gemaaid.
Alle dorpelingen woonden zijn \VUgras{begrafenis} bij. »

%%K t07-c1-13-1 p
13. --- « Ziehier zijn opvolger die op het beekje schaatst. Hij kwam zojuist bij
mij de riem van zijn \VUgras{schaats} \textsuperscript{(26)} laten herstellen
die gebroken was. »

%%K t07-c1-14-1 p
14. --- « Ziehier ook op de oever van het beekje, de nieuwe
\VUgras{veldwachter} \textsuperscript{(27)}. Hij is een oud-gendarme die de
deugnieten van het dorp schrik aanjaagt. Zij vluchten zodra zij zijn
\VUgras{beenkappen} \textsuperscript{(29)} en zijn \VUgras{kepie}
\textsuperscript{(28)} zien. »

%%K t07-c1-15-1 p
15. --- « Ha! ziehier de oude Jozef die een \VUgras{karretje met takkenbossen}
\textsuperscript{(30)} voor de bakker rijdt. --- Dat is waar, zei mijn vader, ik
herken zijn span \VUgras{muildieren} \textsuperscript{(31)}. Ik moet hem
werkelijk spreken, ik zal van de gelegenheid gebruikmaken. Tot ziens, vader
Jacob. »

%%K t07-c1-16-1 p
16. --- Juist toen wij naar buiten gingen, verlieten de kinderen van het dorp de
school en stormden rennend naar de \VUgras{glijbaan} \textsuperscript{(33)} met
gevaar te vallen en bij de \VUgras{val} \textsuperscript{(34)} een lid te
breken. Een van hen nam zijn \VUgras{slee} \textsuperscript{(32)} bij de
wagenmaker, liet een van zijn kameraden erop zitten, en vertrok rennend.

%%K t07-c1-17-1 p
17. --- Anderen stormden naar een \VUgras{sneeuwhoop} \textsuperscript{(37)}
naast de \VUgras{brug} \textsuperscript{(40)} en begonnen het
\VUgras{sneeuwbeeld} \textsuperscript{(39)} te bombarderen dat zij de vorige dag
hadden opgericht en waarop zij een hoed, een pijp en een schooltas over de
schouder gezet hadden.

%%K t07-c1-18-1 p
18. --- Zij bekommeren zich weinig om de ijzige wind en de koude. De meesten
hebben zelfs geen \VUgras{dikke das} \textsuperscript{(35)} en, om zich na het
gevecht met \VUgras{sneeuwballen} \textsuperscript{(38)} weer te warmen, is een
handvol zeer warme \VUgras{kastanjes} \textsuperscript{(42)} hun genoeg, gekocht
bij de oude \VUgras{verkoopster} Jacoba, die van koude beeft onder haar te dunne
\VUgras{versleten sjaal} \textsuperscript{(43)}.

%%K t07-tit-5 sub
\VUsaut{3.0mm}
\VUcentre{12.0pt}{0}{\textit{Tweede tafereel.}}
\VUsaut{3.0mm}

%%K t07-tit-6 sub
\VUpk{10.2pt}{40}{Het landhuis in de lente.}
\VUsaut{4.0mm}

%%K t07-c2-01-1 p
1. --- Ondanks alle genoegens van de winter begroeten wij met diepe vreugde de
terugkeer van de \VUgras{lente}.

%%K t07-c2-01-2 p
Wat een opvrolijkend tafereel biedt een boerenerf, 's avonds, in de maand mei!

%%K t07-c2-02-1 p
2. --- Aan de horizon gaat de \VUgras{zon} \textsuperscript{(1)} langzaam
onder, en overstroomt het \VUgras{land} \textsuperscript{(3)} met haar purperen
\VUgras{stralen} \textsuperscript{(2)}. De vlijtige \VUgras{ploeger}
\textsuperscript{(6)} haast zich zijn \VUgras{akker} \textsuperscript{(4)} te
ploegen.

%%K t07-c2-02-2 p
Met zijn \VUgras{ploeg} \textsuperscript{(5)} voor een krachtig \VUgras{paard}
\textsuperscript{(7)} gespannen, trekt hij de \VUgras{voor}
\textsuperscript{(8)} waarin de \VUgras{zaaier} \textsuperscript{(9)} met volle
hand het \VUgras{zaad} zal werpen dat de gouden aren zal opleveren.

%%K t07-c2-03-1 p
3. --- De \VUgras{maaiers} \textsuperscript{(11)} met hun zeisen hebben het gras
van de weide gemaaid, de hooiers hebben het \VUgras{hooi}
\textsuperscript{(10)} gedroogd door het met \VUgras{vorken}
\textsuperscript{(12)} en harken te keren, en ziehier dat zij terugkomen.

%%K t07-c2-03-2 p
Sommigen lopen voor de zware wagen die door ossen getrokken wordt; zij dragen
hun gereedschap over de schouder. Anderen, luier, hebben zich behaaglijk op het
geurende hooi geïnstalleerd.

%%K t07-c2-04-1 p
4. --- In het voorbijgaan groeten zij vriendschappelijk de \VUgras{tuinman}
\textsuperscript{(16)} die in de \VUgras{boomgaard} \textsuperscript{(13)} bezig
is de \VUgras{vruchtbomen} te snoeien en de \VUgras{perenbomen}
\textsuperscript{(14)} te enten. De \VUgras{pruimenbomen}
\textsuperscript{(15)} beginnen te bloeien, en, in de goed bemeste
\VUgras{moestuin} \textsuperscript{(17)}, staan de groenten, kolen en sla er al
goed bij.

%%K t07-c2-04-2 p
Op het muurtje spreidt een mooie \VUgras{pauw} \textsuperscript{(18)} zijn
staart uit, om zijn prachtige veren te laten bewonderen.

%%K t07-tit-7 sub
\VUpk{10.2pt}{40}{Boerenerf.}
\VUsaut{3.0mm}

%%K t07-c2-05-1 p
5. --- De deuren van de \VUgras{paardenstal} \textsuperscript{(19)} staan open
en men ziet de ruif en het \VUgras{strobed} \textsuperscript{(23)} van de
\VUgras{merrie} \textsuperscript{(21)} die de \VUgras{stalknecht}
\textsuperscript{(20)} zojuist heeft uitgespannen. Het \VUgras{veulen}
\textsuperscript{(22)} zo grappig met zijn lange benen, volgt zijn moeder
huppelend.

%%K t07-c2-06-1 p
6. --- Bij de \VUgras{put} \textsuperscript{(29)} gaan de \VUgras{koeien}
\textsuperscript{(25)} drinken uit de houten \VUgras{trog}
\textsuperscript{(26)} die hun als drinkbak dient. Het \VUgras{kalf}
\textsuperscript{(24)} maakt van de gelegenheid gebruik om te zuigen, en de
\VUgras{stier} \textsuperscript{(27)} die de goede geur van de hooiberg ruikt,
loeit vrolijk en heft trots de kop met krachtige \VUgras{horens}
\textsuperscript{(28)}.

%%K t07-c2-07-1 p
7. --- Allen werken. De \VUgras{dienstmeid} \textsuperscript{(32)} houdt het
\VUgras{touw} \textsuperscript{(31)} van de put in de hand. Zij put water met
een \VUgras{emmer} \textsuperscript{(30)} om het vee te drenken. Naast de
\VUgras{schuur} \textsuperscript{(33)} harkt een man het stro bijeen en voor de
deur van het \VUgras{woonhuis} \textsuperscript{(34)} spint een oude
\VUgras{spinster} \textsuperscript{(35)} met haar spinnewiel en haar
\VUgras{spinrok}, het \VUgras{vlas} of de \VUgras{hennep} zoals in vroeger
tijden.

%%K t07-tit-8 sub
\VUsaut{3.0mm}
\VUpk{10.2pt}{40}{De boer.}
\VUsaut{3.0mm}

%%K t07-c2-08-1 p
8. --- De \VUgras{boer} \textsuperscript{(36)} leidt al die werkzaamheden en
geeft aanwijzingen aan zijn knechten. Hij beveelt de \VUgras{eg}
\textsuperscript{(40)} en het \VUgras{houweel} \textsuperscript{(39)} onder dak
te brengen die men naast het \VUgras{hok} \textsuperscript{(38)} van de
\VUgras{hond} \textsuperscript{(37)} vergeten heeft.

%%K t07-c2-09-1 p
9. --- Zijn geroep jaagt een \VUgras{duif} \textsuperscript{(41)} schrik aan,
die met alle vleugels de \VUgras{duiventil} \textsuperscript{(42)} in vliegt, en
de \VUgras{hennen} \textsuperscript{(43)} die zich haasten het \VUgras{kippenhok}
\textsuperscript{(44)} weer binnen te gaan.

%%K t07-c2-10-1 p
10. --- Voor de \VUgras{schaapskooi} \textsuperscript{(48)}, ziehier de oude
\VUgras{herder} \textsuperscript{(45)} die zijn \VUgras{(schapen)kudde}
\textsuperscript{(49)} terugbrengt. Met zijn \VUgras{herdersstaf}
\textsuperscript{(46)} in de hand, die hem nooit ontbreekt, evenmin als zijn
verschoten \VUgras{kiel} \textsuperscript{(47)}, drijft hij naar de stal
\VUgras{rammen} \textsuperscript{(50)}, \VUgras{ooien}
\textsuperscript{(53)}, \VUgras{lammeren} \textsuperscript{(54)},
\VUgras{geiten} \textsuperscript{(51)} en \VUgras{geitjes}
\textsuperscript{(52)}. Zijn trouwe \VUgras{hond} \textsuperscript{(55)} Argus
komt als laatste en spoort met zijn geblaf de achterblijvers tot spoed aan.

%%K t07-c2-11-1 p
11. --- Al dat lawaai en die drukte verstoren de rust niet van de goede
\VUgras{melkkoe} \textsuperscript{(56)} die haar \VUgras{belletje}
\textsuperscript{(57)} laat rinkelen, terwijl men haar voor de deur van de
\VUgras{stal} \textsuperscript{(58)} melkt.

%%K t07-c2-12-1 p
12. Zij schijnt te begrijpen, dat zij haar melk moet geven, opdat de
\VUgras{boerin} \textsuperscript{(60)} \VUgras{boter} in de \VUgras{karn}
\textsuperscript{(61)} zou kunnen maken of de uitstekende \VUgras{kazen}
\textsuperscript{(64)}, die van de plank druipen die op de \VUgras{tobbe}
\textsuperscript{(63)} ligt.

%%K t07-c2-13-1 p
13. --- De hele \VUgras{melkerij} \textsuperscript{(59)} wordt zeer schoon
gehouden door de \VUgras{boerenknecht} \textsuperscript{(65)} die zojuist de
\VUgras{melkkannen} \textsuperscript{(62)} gewassen heeft in de grote emmer die
hij in de hand draagt.

%%K t07-tit-9 sub
\VUsaut{3.0mm}
\VUpk{10.2pt}{40}{Het pluimvee.}
\VUsaut{3.0mm}

%%K t07-c2-14-1 p
14. --- Met zijn dikke klompen loopt hij wat log, en zo jaagt hij de
\VUgras{kalkoen} \textsuperscript{(68)}, de \VUgras{eenden}
\textsuperscript{(66)}, de \VUgras{woerden} \textsuperscript{(67)} en de
\VUgras{ganzen} \textsuperscript{(69)} op de vlucht die \VUgras{hun snavel}
\textsuperscript{(70)} wijd opensperren en onrustig schreeuwen.

%%K t07-c2-14-2 p
De \VUgras{haan} \textsuperscript{(71)}, strijdlustiger, richt zijn rode
\VUgras{kam} \textsuperscript{(72)} op, schudt de veren van zijn \VUgras{staart}
\textsuperscript{(73)} en laat een luid « kukeleku » horen.

%%K t07-c2-15-1 p
15. --- Een \VUgras{moederkloek} \textsuperscript{(74)} scharrelt op de
\VUgras{mesthoop} \textsuperscript{(81)} met haar \VUgras{kuikentjes}
\textsuperscript{(75)}. Het \VUgras{biggetje} \textsuperscript{(78)} vlucht naar zijn
moeder, de enorme \VUgras{zeug} \textsuperscript{(77)} die wij naast het
varkenskot zien.

%%K t07-c2-16-1 p
16. --- In hun hokken eten de \VUgras{konijnen} \textsuperscript{(79)} gretig het
malse \VUgras{gras} \textsuperscript{(80)} dat de dochter van de boer hun brengt.
Jantje houdt veel van haar konijnen en, elke dag, komt zij haar vriendjes
bezoeken nadat zij de verse \VUgras{eieren} \textsuperscript{(76)} in het
kippenhok gehaald heeft.
\VUsaut{6.0mm}
\VUfilet{22mm}
